\documentclass[11pt,reqno]{amsart}

\usepackage[T1]{fontenc}
\usepackage{lmodern}
\usepackage{amsmath,amssymb,mathtools}
\usepackage{booktabs}
\usepackage{enumitem}
\usepackage{microtype}
\usepackage[colorlinks=true,linkcolor=blue,citecolor=blue,urlcolor=blue]{hyperref}

\hypersetup{
  pdftitle={Endpoint Localization for Type III and Type IV Salem--Boyd Interlacing Families},
  pdfauthor={KEO},
  pdfsubject={Salem and Pisot numbers, interlacing, and endpoint compactification}
}

\newtheorem{theorem}{Theorem}[section]
\newtheorem{proposition}[theorem]{Proposition}
\newtheorem{lemma}[theorem]{Lemma}
\newtheorem{corollary}[theorem]{Corollary}
\theoremstyle{definition}
\newtheorem{definition}[theorem]{Definition}
\newtheorem{example}[theorem]{Example}
\theoremstyle{remark}
\newtheorem{remark}[theorem]{Remark}

\DeclareMathOperator{\ord}{ord}
\DeclareMathOperator{\Res}{Res}
\DeclareMathOperator{\TV}{TV}
\newcommand{\D}{\mathbb D}
\newcommand{\T}{\mathbb T}
\newcommand{\Q}{\mathbf Q}
\newcommand{\Z}{\mathbf Z}
\newcommand{\R}{\mathbf R}
\newcommand{\calS}{\mathcal S}
\newcommand{\calT}{\mathcal T}

\title[Salem--Boyd endpoint localization]{Endpoint Localization for Type III
and Type IV\\Salem--Boyd Interlacing Families}
\author{KEO}
\date{September 2026}
\keywords{Salem numbers, Pisot numbers, interlacing, Herglotz functions,
endpoint compactification}

\begin{document}

\begin{abstract}
McKee and Smyth proved that every Salem number admits at least one, generally
non-unique, interlacing witness of Types I--IV.  We study relative sequential
closure in $(1,\infty)$, with emphasis on the moving-pole Types III and IV.
Type I is closed in the interior, while an exact CS-to-CC transformation puts
Type II in a closed CC chart.

For Type III we prove closure away from collision between the output and the
denominator's exterior root, as well as on the positive-residue and
bounded-denominator-degree strata.  Every remaining collision has
unbounded-degree Salem denominator roots, and its residue is strictly
linearly comparable with the pole--output gap; in particular, a higher-order
sublinear collision is impossible.  We also give an explicit Type III
witness showing that the Type II Weyl descent cannot be iterated inside any
fixed integral Laurent lattice.

For Type IV an exact nonnegative defect identity yields the universal bound
$c<1/2$, a quantitative pole--output separation, and a three-stratum
subsequential compactification: finite separation, zero-residue collision,
or half-residue escape.  The finite stratum, every compact residue window in
$(0,1/2)$, and every sublinear collision are closed.  Half-residue escape has
an explicit coefficient-growth fingerprint.  Consequently every
hypothetical bad limit is confined, after parameter-convergent extraction,
to a strict first-order Type III or Type IV collision, or to Type IV
half-residue escape.  A witness-independent normal height must simultaneously
tend to infinity, while remaining at most linear in half the Salem degree.
These results prove neither Boyd's global closedness conjecture nor the
stronger derived-set identity; in particular they do not treat accumulation
at $1$.
\end{abstract}

\maketitle

\section{Introduction and scope}

A Salem number is a real algebraic integer $\tau>1$ whose conjugates consist
of $\tau^{-1}$ and points on the unit circle, with at least one unit-circle
conjugate.  A Pisot number is a real algebraic integer larger than one whose
other conjugates lie in the open unit disk.  We write $\calT$ and $\calS$ for
the respective sets.

Boyd formulated two related conjectures; see \cite{boyd77,boyd78}.  The
closedness conjecture asserts that $\calS\cup\calT$ is closed in $\R$.  The
interior sequential component studied here is the inclusion
$\calT'\cap(1,\infty)\subseteq\calS\cup\calT$.
Accumulation at $1$ is a separate Salem-gap problem, implied in particular by
Lehmer's conjecture.  Boyd's stronger derived-set conjecture is
$(\calS\cup\calT)'=\calS$.  We prove neither global conjecture.  We instead
localize every possible failure of the displayed interior component within a
specified interlacing-witness subsequence.

During revision we became aware of Amara \cite{amara25}, whose 2025 preprint
claims that every accumulation point of $\calT$ belongs to $\calS$, together
with related closedness and Lehmer conclusions.  As of 4 September 2026 the
arXiv record lists only its initial version and no journal reference.  We do
not use or assess that manuscript, and nothing below depends on the status of
its announced conclusions.  The present paper is an independent quantitative
compactification and endpoint-localization result, not a claim to settle
either Boyd conjecture.

McKee and Smyth's Theorem~7.3 \cite{mckeesmyth} gives every Salem number at
least one witness of Types I--IV; the witness need not be unique.  Types I
and II are governed by CC and CS interlacing, whereas Types III and IV are
the type-1 and type-2 SS cases.  Types III and IV contain moving exterior
poles, yet possess rigid compactifications whose possible failures occur only
at explicit endpoint strata.  The emphasis on Type IV is also natural in
view of \cite[Thm.~9.2]{mckeesmyth}: if a least Salem number exists, then it
is of Type IV and of no other type.  McKee and Smyth already identified Pisot
limits for specific limit-function constructions \cite[Thms.~4.2 and
5.3]{mckeesmyth}; our concern is different, namely sequences for which the
witness itself varies.

Our arguments combine three simple ideas.  First, after reversing the
interlacing quotient, the Taylor series at the origin has integer
coefficients.  Second, Herglotz representations yield normal families and
uniform coefficient bounds.  Third, Cantor's rationality theorem for
integer power series of bounded characteristic \cite{cantor} turns a
coefficientwise limit back into a rational function.  Meromorphic Hurwitz
then transports the unique output zero unless it collides with a pole.

The main results are Theorems~\ref{thm:typeIII}, \ref{thm:typeIVcompact},
\ref{thm:IVsublinear}, and \ref{thm:aggregate}.  Theorem~\ref{thm:typeII}
removes Type II by an exact algebraic reduction rather than an asymptotic
argument.  Theorem~\ref{thm:heightclosure} provides a complementary closure
principle independent of the chosen McKee--Smyth witness.

For reference, the logical output of the paper is summarized as follows.
\begin{center}
\footnotesize
\begin{tabular}{@{}p{0.18\textwidth}p{0.71\textwidth}@{}}
\toprule
chart & proved closure or remaining degeneration\\
\midrule
Type I & interior closure\\
Type II & exact reduction to a closed CC output equation\\
Type III & closed off collision and in bounded degree; otherwise an
unbounded-degree Salem twin with residue comparable to the gap\\
Type IV & finite-separation and interior-residue strata closed; only
first-order collision or half-residue escape remains\\
normal height & bounded height is closed; every bad sequence has
$\kappa\to\infty$ and $\kappa=O(\deg R)$\\
\bottomrule
\end{tabular}
\end{center}

\section{Analytic and arithmetic compactness}

We retain the terminology and endpoint conventions of
\cite[Secs.~2.2--2.5 and Thm.~7.3]{mckeesmyth}.  Thus CC means
cyclotomic--cyclotomic interlacing, CS means cyclotomic--Salem interlacing,
and SS means Salem--Salem interlacing, with the usual caveat that the
interlacing polynomials need not themselves be cyclotomic or minimal Salem
polynomials.  A Type I, II, III, or IV witness uses, respectively, a CC, CS,
type-1 SS, or type-2 SS quotient $Q/P$, where $P,Q\in\Z[z]$ are coprime,
monic, of the same degree, and
\begin{equation}\label{eq:witness}
 \frac{Q(z)}{(z-1)P(z)}=\frac2{z+1}
\end{equation}
has as its nonzero solutions the output Salem orbit and possibly roots of
unity.  For $\nu\in\{\mathrm I,\mathrm{II},\mathrm{III},\mathrm{IV}\}$,
let $\calT_\nu$ be the Salem outputs admitting a witness of Type $\nu$.
These classes need not be disjoint.

For a polynomial $F$ of degree $d$, put
\[
 F^*(t)=t^dF(t^{-1}).
\]
The reversed quotient associated with an interlacing witness is written
\begin{equation}\label{eq:Phi}
 \Phi(t)=\frac{Q^*(t)}{P^*(t)}.
\end{equation}
With monic integral input, $P^*(0)=Q^*(0)=1$, so $\Phi\in\Z[[t]]$.
For the original Type I, III, and IV charts, the output equation has the form
\begin{equation}\label{eq:output}
 G(t):=(1+t)\Phi(t)-2(1-t)=0,
\end{equation}
and the Salem output corresponds to the unique zero $t=\tau^{-1}$ in $\D$.
The CC chart obtained from Type II has a different equation, recorded in
Theorem~\ref{thm:CC31} below.

We use the following input from the interlacing classification.  The endpoint
alternatives for CS pairs, the real Cauchy transforms, and the signs of their
residues are those of \cite[Secs.~2.3--2.5 and Thm.~7.3]{mckeesmyth}.
For a monic CC pair, the reversed quotient is a Carath\'eodory function:
\begin{equation}\label{eq:caratheodory}
 \Phi(0)=1,\qquad \Re\Phi(t)>0\quad(t\in\D),\qquad
 |[t^n]\Phi|\le2\quad(n\ge1).
\end{equation}
Indeed the real partial fractions have positive residues on $[-2,2]$;
lifting their atoms to $\T$ gives the Herglotz representation
\[
 \Phi(t)=\int_{\T}\frac{\zeta+t}{\zeta-t}\,d\mu(\zeta),
 \qquad \mu(\T)=1.
\]
For Type III the same partial-fraction description has one positive exterior
reciprocal pair, with the type-1 SS ordering $1<\rho<\tau$.  For Type IV it has
exactly two negative exterior residues and a positive symmetric compact
part.  These facts supply the decompositions used in
\eqref{eq:typeIIIdecomp} and \eqref{eq:typeIVreal}; all estimates after those
decompositions are proved here.

We use the following standard arithmetic form of Cantor's theorem.

\begin{theorem}[Cantor's bounded-characteristic theorem]\label{thm:cantor}
Let $F$ be holomorphic in a neighborhood of $0$, with Taylor coefficients in
$\Z$.  If this germ admits a meromorphic continuation of bounded
characteristic to $\D$, then $F\in\Q(t)$.  In reduced form one may write
\[
 F=A/B,\qquad A,B\in\Z[t],\qquad B(0)=1.
\]
\end{theorem}

\begin{proof}
The rationality assertion is Cantor's theorem \cite{cantor}.  For the
integral normalization, write $F=\sum a_nt^n$ and consider the Hankel columns
$v_n=(a_n,a_{n+1},\ldots)$.  Rationality makes their $\Q$-span finite
dimensional.  A finite coordinate projection is injective on that span, so
the subgroup generated by the $v_n$ embeds in some $\Z^r$ and is finitely
generated.  The shift $v_n\mapsto v_{n+1}$ is an endomorphism of this
lattice.  Cayley--Hamilton therefore gives a monic integral recurrence for
the $a_n$.  Its generating-function denominator lies in $\Z[t]$ and has
constant term one; the numerator is then integral as well.
\end{proof}

\begin{lemma}[Integer-germ compactness]\label{lem:compactness}
Suppose $\Phi_j\in\Z[[t]]$ are meromorphic on $\D$.  Assume that there is
$\delta>0$ such that every $\Phi_j$ is holomorphic on $|t|<\delta$, the family
is normal there, and all moving poles lie outside that disk.  Assume also
that every fixed Taylor coefficient is bounded and that a subsequential
meromorphic limit on pole-free compacta has bounded characteristic.  Then a
further subsequence converges coefficientwise to an integer series $\Phi$,
the meromorphic limit represents that series near zero, and
$\Phi\in\Q(t)$.
\end{lemma}

\begin{proof}
Diagonal extraction makes every integer coefficient eventually constant.
Local uniform convergence on $|t|<\delta$ identifies those stabilized
coefficients with the derivatives of the meromorphic limit at zero.  The
limit germ therefore has integer coefficients, and
Theorem~\ref{thm:cantor} applies.
\end{proof}

If $\Phi=A/B$ is reduced, then at a pole $p\in\D$ the two cleared numerators
used below satisfy
\begin{equation}\label{eq:coprimepole}
 \begin{aligned}
 H_{\mathrm I}(p)&=(1+p)A(p)\ne0,\\
 H_{3.1}(p)&=2pA(p)\ne0,
 \end{aligned}
\end{equation}
where
\[
 H_{\mathrm I}=(1+t)A-2(1-t)B,\qquad
 H_{3.1}=2tA-(1-t^2)B.
\]
Here $p\ne-1$, $p\ne0$ because $B(0)=1$, and $(A,B)=1$.  Thus zeros cannot be
hidden at denominator poles.  We apply meromorphic Hurwitz only on
exhaustions avoiding those poles, together with explicit output--pole
separation.

\begin{lemma}[Integral reversal criterion]\label{lem:reversal}
Let $H\in\Z[t]$ satisfy $H(0)=-1$.  If $H$ has exactly one zero
$u\in(0,1)$ in $\D$, counted with multiplicity, then
$u^{-1}\in\calS\cup\calT$.
\end{lemma}

\begin{proof}
If $m=\deg H$, then
$-z^mH(z^{-1})$ is monic integral, has $u^{-1}>1$ as a root, and all its
other roots in the closed unit disk.  The irreducible factor containing
$u^{-1}$ is therefore a Pisot or Salem minimal polynomial.
\end{proof}

\begin{proposition}[Type I closure]\label{prop:typeI}
The Type I output class is closed in the interior in the sense that
\[
 \overline{\calT_{\mathrm I}}\cap(1,\infty)
 \subseteq\calS\cup\calT.
\]
\end{proposition}

\begin{proof}
Let $\tau_j\to\theta>1$ and choose monic CC witnesses.  By
\eqref{eq:caratheodory}, their integer reversed quotients $\Phi_j$ form a
normal family and every coefficient belongs to a fixed finite set.  A
subsequence converges locally uniformly to an integer Carath\'eodory germ
$\Phi$, which is rational by Theorem~\ref{thm:cantor}; write it as $A/B$ in
the integral normalized form.  The functions in \eqref{eq:output} have the
unique disk zeros $u_j=\tau_j^{-1}\to\theta^{-1}$.  Hurwitz, applied on every
compact subdisk, shows that the limit has exactly that one disk zero.  The
first identity in \eqref{eq:coprimepole} shows that the cleared polynomial
$H_{\mathrm I}$ has the same disk zeros and $H_{\mathrm I}(0)=-1$.
Lemma~\ref{lem:reversal} finishes the proof.
\end{proof}

\begin{theorem}[Interior closure for the CC output equation]\label{thm:CC31}
Let $U_j$ be monic integral polynomials and $C_j$ integral polynomials of the
same degree such that $(2C_j)/U_j$ is a positive CC-interlacing quotient.  Put
\[
 \Psi_j=C_j^*/U_j^*,\qquad k_j=\Psi_j(0)\in\Z_{>0}.
\]
Suppose $u_j=\tau_j^{-1}$ is the unique disk solution of
\begin{equation}\label{eq:CC31equation}
 2t\Psi_j(t)=1-t^2
\end{equation}
and $\tau_j\to\theta>1$.  Then $\theta\in\calS\cup\calT$.
\end{theorem}

\begin{proof}
The McKee--Smyth numerator is $2C_j$, while $\Psi_j$ is its half-normalized
reversed quotient.  Thus $\Psi_j$ is Carath\'eodory of mass $k_j$.
Harnack's inequality
on the positive radius and \eqref{eq:CC31equation} give
\[
 k_j\frac{1-u_j}{1+u_j}\le\Psi_j(u_j)
 =\frac{1-u_j^2}{2u_j},
 \qquad
 k_j\le\frac{(1+u_j)^2}{2u_j}.
\]
Thus the positive integer masses are bounded, and every fixed Taylor
coefficient is bounded by twice that mass.  Herglotz compactness and
Theorem~\ref{thm:cantor} give a rational integer-germ limit
$\Psi=A/B$.  Hurwitz preserves the unique zero $\theta^{-1}$.  By the second
identity in \eqref{eq:coprimepole}, the cleared polynomial
$H_{3.1}=2tA-(1-t^2)B$ has exactly that disk zero and satisfies
$H_{3.1}(0)=-1$.  Apply Lemma~\ref{lem:reversal}.
\end{proof}

\section{The exact Type II reduction}

We include Type II because it removes an apparent endpoint family before the
Type III/IV analysis begins.  Let $P,Q\in\Z[z]$ be a monic CS-interlacing pair
of the same degree producing a Type II output.  Thus the nonunit solutions of
\begin{equation}\label{eq:typeIIeq}
 \frac{Q(z)}{(z-1)P(z)}=\frac2{z+1}
\end{equation}
are the prescribed Salem root, its conjugates, and allowed unit roots.

\begin{theorem}[CS-to-CC Weyl complement]\label{thm:typeII}
Put
\begin{equation}\label{eq:UC}
 U=\frac{Q}{z-1},\qquad
 C=\frac{(z+1)Q-(z-1)P}{z}.
\end{equation}
By the CS endpoint alternatives in \cite[Sec.~2.3]{mckeesmyth}, if
$r=\ord_{z=1}Q$, then $r\in\{1,3\}$.  Define
\[
 g=\begin{cases}1,&r=1,\\z-1,&r=3,\end{cases}\qquad
 \widetilde U=U/g,\quad \widetilde C=C/g.
\]
Then $2\widetilde C/\widetilde U$ is an admissible CC quotient for the
McKee--Smyth Theorem~3.1 construction, and
\begin{equation}\label{eq:typeIIidentity}
 2z\widetilde C-(z^2-1)\widetilde U
 =\frac{(z+1)Q-2(z-1)P}{g}.
\end{equation}
Thus at most the cyclotomic factor $z-1$ is removed, every nonunit output root
is preserved, and
\[
 \calT_{\mathrm{II}}\subseteq\calT_{3.1}.
\]
Here $\calT_{3.1}$ denotes the output class of the CC construction in
\cite[Thm.~3.1]{mckeesmyth}.
In particular,
\begin{equation}\label{eq:typeIIclosure}
 \overline{\calT_{\mathrm{II}}}\cap(1,\infty)
 \subseteq\calS\cup\calT.
\end{equation}
\end{theorem}

\begin{proof}
The CS endpoint convention makes the multiplicity of $1$ in $Q$ odd and at
most three.  Since $P(0)=1$ and $Q(0)=-1$, the second numerator in
\eqref{eq:UC} is divisible by $z$; its leading terms cancel, so $C$ is
integral of degree at most $\deg Q-1$, while $U$ is monic.  Reciprocity shows
that $U$ is reciprocal and $C$ antireciprocal.

In the real-line variable, write the CS quotient as a positive-residue
Cauchy transform $m(x)\sim x^{-1}$.  Its reciprocal Weyl complement
\[
 L(x)=x-\frac1{m(x)}
\]
has positive residue $-1/m'(\beta)$ at every zero $\beta$ of $m$, all of
which lie in $[-2,2]$, and $L=O(x^{-1})$.  Hence it is again a positive
Cauchy transform with circularly interlacing numerator and denominator.
Its total mass
\[
 k=\lim_{x\to\infty}xL(x)>0
\]
is the leading coefficient of $C$, so $k\in\Z_{>0}$ and
$\deg C=\deg U$.
Direct substitution identifies it with the transform of $C/U$.  Scaling the
numerator by $2$ does not change its zeros or the CC interlacing, so
$2C/U$ is a CC quotient as well.  The only common factor is $z-1$, absent
for $r=1$ and simple after cancellation for $r=3$.  Clearing the transformed equation gives
\eqref{eq:typeIIidentity}.

For completeness, let $a>2$ be the largest CS pole.  The CS ordering gives
$m(x)<0$ on $(2,a)$, so
\[
 L(2+)=2-\frac1{m(2+)}>2
\]
(with value $+\infty$ at an endpoint zero).  Thus the real transform of
$2\widetilde C/\widetilde U$ satisfies the endpoint condition in
\cite[Thm.~3.1]{mckeesmyth}.  If
$\Psi=\widetilde C^*/\widetilde U^*$, reversing
\eqref{eq:typeIIidentity} at a nonunit root gives exactly
\[
 2t\Psi(t)=1-t^2.
\]
This explains both the factor $2$ in the McKee--Smyth numerator and the
half-normalized choice of $\Psi$ in Theorem~\ref{thm:CC31}.
Theorem~\ref{thm:CC31} therefore proves the asserted interior closure.
\end{proof}

\section{Why the Type II descent does not iterate}

The Weyl complement in Theorem~\ref{thm:typeII} is geometrically tempting in
Type III as well.  The next proposition shows that its arithmetic
normalization can fail at every depth; this is why the direct integer-germ
compactness argument in the following section is needed.

\begin{proposition}[Failure of fixed-lattice Weyl iteration]
\label{prop:weylfailure}
There is a monic Type III witness whose first Weyl complement is again a
positive-residue Cauchy transform, but whose second normalized complement has
Laurent coefficients with unbounded powers of $2$ in their denominators.
Consequently no fixed nonzero rational scaling puts all of those coefficients
in $\Z$.
\end{proposition}

\begin{proof}
Consider
\begin{align*}
 P(z)&=z^6-2z^5-z^4-5z^3-z^2-2z+1,\\
 Q(z)&=z^6-2z^5+2z-1.
\end{align*}
Its real interlacing quotient is
\begin{equation}\label{eq:weylm}
 m(x)=\frac{x(x^4-6x^2+7)}
 {(x^2-1)(x^2-x-3)(x^2+x-3)}.
\end{equation}
The numerator zeros are
\[
 -\sqrt{3+\sqrt2},\ -\sqrt{3-\sqrt2},\ 0,\
 \sqrt{3-\sqrt2},\ \sqrt{3+\sqrt2},
\]
and they strictly alternate with the six displayed denominator zeros.  The
largest denominator zero $(1+\sqrt{13})/2$ is larger than the largest
numerator zero, so this is the type-1 SS ordering.  Moreover
\begin{equation}\label{eq:weylfactor}
 (z+1)Q-2(z-1)P
 =-(z-1)(z^2+1)(z^4-4z^3-z^2-4z+1).
\end{equation}
The last factor has trace polynomial $X^2-4X-3$ and is irreducible by
Eisenstein at $7$ after replacing $z$ by $z+1$.  Thus the pair is a genuine
Type III witness.

The first Weyl complement and its total mass are
\begin{equation}\label{eq:firstweyl}
 L(x)=x-\frac1{m(x)}
 =\frac{(x^2-3)(2x^2-3)}{x(x^4-6x^2+7)},
 \qquad \lim_{x\to\infty}xL(x)=2.
\end{equation}
The positive normalized next step is therefore forced to be
\begin{equation}\label{eq:secondweyl}
 M(x)=x-\frac2{L(x)}
 =\frac{x(3x^2-5)}{(x^2-3)(2x^2-3)}.
\end{equation}
Put $t=x^{-2}$ and write $M=x^{-1}\sum_{n\ge0}a_nt^n$.  Directly from
\eqref{eq:secondweyl},
\[
 (2-9t+9t^2)\sum_{n\ge0}a_nt^n=3-5t.
\]
If $b_n=2^{n+1}a_n$, then $b_0=3$, $b_1=17$, and
\[
 b_n=9b_{n-1}-18b_{n-2}\qquad(n\ge2).
\]
Every $b_n$ is odd, so $v_2(a_n)=-(n+1)$.  For a fixed
$q\in\Q^\times$, the valuations $v_2(qa_n)$ are therefore negative for all
sufficiently large $n$.  No fixed scaling restores an integral Laurent
lattice.
\end{proof}

\section{Type III: noncollision closure}

Fix one Type III witness for each output $\tau_j$.  In the disk its reversed
quotient has one exterior pole pair.  Let $1<\rho_j<\tau_j$ be the exterior
root of the denominator and let $\lambda_j>0$ be the residue assigned to
each member of that reciprocal pair.  Put $u_j=\tau_j^{-1}$.

\begin{proposition}[Type III decomposition and sharpened budgets]
\label{prop:IIIbudget}
There is a real-symmetric Carath\'eodory function $C_j$ such that
\begin{equation}\label{eq:typeIIIdecomp}
 \Phi_j(t)=C_j(t)+
 \frac{2\lambda_j(1-t^2)}
 {(1-\rho_jt)(1-\rho_j^{-1}t)},\qquad
 C_j(0)=1-2\lambda_j\ge0.
\end{equation}
At the output zero define
\[
 \delta_j=\frac{(1+u_j)C_j(u_j)}{2(1-u_j)}.
\]
Then $0\le\delta_j<1$ and the exact defect factorization is
\begin{equation}\label{eq:typeIIIdefect}
 \lambda_j=
 \frac{(\tau_j-\rho_j)(\rho_j\tau_j-1)}
 {\rho_j(\tau_j+1)^2}(1-\delta_j).
\end{equation}
In particular both bounds
\begin{align}
 \lambda_j&\le
 \frac{(\tau_j-\rho_j)(\rho_j\tau_j-1)}
 {\rho_j(\tau_j+1)^2},\label{eq:typeIIIbudget}\\
 \lambda_j&\le
 \frac{(\tau_j-\rho_j)(\rho_j\tau_j-1)}
 {2\tau_j(\rho_j+1)^2}\label{eq:typeIIIharnackbudget}
\end{align}
hold globally.
\end{proposition}

\begin{proof}
Remove the two positive exterior atoms from the positive Cauchy measure of
the Type III quotient; the remaining symmetric atoms lie on $[-2,2]$ and
give $C_j$.
Their masses add to $1-2\lambda_j$, proving \eqref{eq:typeIIIdecomp}.
Put
\[
 D_j(u)=(1-\rho_ju)(1-\rho_j^{-1}u),\qquad
 h(u)=\frac{2(1-u)}{1+u}.
\]
The exact output equation is
\begin{equation}\label{eq:typeIIIdefect0}
 h(u_j)-C_j(u_j)=
 \frac{2\lambda_j(1-u_j^2)}{D_j(u_j)}>0.
\end{equation}
Solving this identity and substituting $u_j=\tau_j^{-1}$ gives
\eqref{eq:typeIIIdefect}; it also shows $0\le\delta_j<1$ and gives
\eqref{eq:typeIIIbudget}.

Harnack's inequality on the positive radius gives
\[
 C_j(u_j)\ge C_j(0)\frac{1-u_j}{1+u_j}
 =(1-2\lambda_j)\frac{1-u_j}{1+u_j}.
\]
Combining this with \eqref{eq:typeIIIdefect0} yields
\[
 2\lambda_j\bigl((1+u_j)^2-D_j(u_j)\bigr)\le D_j(u_j).
\]
Since
\[
 (1+u_j)^2-D_j(u_j)=\frac{(\rho_j+1)^2}{\rho_j\tau_j},
\]
the second bound follows.
\end{proof}

\begin{lemma}[Simplicity of the Type III output]\label{lem:IIIsimple}
The disk output zero of a Type III witness is simple.
\end{lemma}

\begin{proof}
In the real variable the positive-residue odd Cauchy transform has the form
\[
 m(x)=\frac{\eta_0}{x}+
 \sum_{a>0}\eta_a\left(\frac1{x-a}+\frac1{x+a}\right),
 \qquad \eta_a>0,
\]
and the output equation is $xm(x)=2$.  To the right of the largest pole,
\[
 \frac d{dx}\bigl(xm(x)\bigr)
 =-4x\sum_{a>0}\frac{\eta_a a^2}{(x^2-a^2)^2}<0,
\]
where strictness follows from the exterior pole.  The disk change of variable
is locally nonsingular at the output, proving the claim.
\end{proof}

\begin{theorem}[Type III noncollision closure]\label{thm:typeIII}
Let $\tau_j\to\theta>1$ be Type III Salem outputs with chosen witnesses.
If
\begin{equation}\label{eq:typeIIIseparation}
 \liminf_j(\tau_j-\rho_j)>0,
\end{equation}
then $\theta\in\calS\cup\calT$.  Within the compactification
\eqref{eq:typeIIIdecomp}, the only way for the output zero to be lost is
\begin{equation}\label{eq:typeIIIcollision}
 \tau_j-\rho_j\longrightarrow0,\qquad
 \lambda_j\longrightarrow0.
\end{equation}
Every $\rho_j$ is a Salem number or a reciprocal-quadratic Pisot number.
\end{theorem}

\begin{proof}
After a subsequence, $\rho_j\to\rho\ge1$, $\lambda_j\to\lambda$, and
$C_j\to C$ locally uniformly.  Each fixed Taylor coefficient of
\eqref{eq:typeIIIdecomp} is bounded uniformly in $j$.  Moreover
$\rho_j<\tau_j=O(1)$, so the disk poles $\rho_j^{-1}$ stay outside a fixed
neighborhood $|t|<\delta$ of zero.  The limit is a Carath\'eodory function
plus one rational pole pair and hence has bounded characteristic.  Since the
full germs are integral, Lemma~\ref{lem:compactness} gives a rational limit
$\Phi=A/B$ in the normalized integral form of
Theorem~\ref{thm:cantor}.

Under \eqref{eq:typeIIIseparation}, $u_j$ remains separated from the moving
pole $\rho_j^{-1}$.  Choose a closed disk about $u=\theta^{-1}$ disjoint
from the limiting pole set.  For large $j$, $G_j$ is holomorphic there and
has its one simple zero $u_j$ there, so the argument principle gives one
limit zero counted with multiplicity.  A zero elsewhere on a pole-free
compact disk would, by Hurwitz, force a second disk zero of $G_j$.  The first
identity in \eqref{eq:coprimepole} excludes a zero created by clearing a
reduced denominator pole.  Hence $H_{\mathrm I}$ has exactly the one disk
zero $u$, and Lemma~\ref{lem:reversal} gives
$\theta\in\calS\cup\calT$.

It remains to identify the only loss mode.  If $\rho>1$ and $\lambda=0$, put
\[
 D_j=(1-\rho_jt)(1-\rho_j^{-1}t),\qquad K_j=D_jG_j.
\]
Then $K_j$ is holomorphic, its only disk zero is $u_j$, and
\[
 K_j(\rho_j^{-1})=
 2\lambda_j(1+\rho_j^{-1})(1-\rho_j^{-2})\ne0.
\]
But
\[
 K_j\longrightarrow
 (1-\rho t)(1-\rho^{-1}t)\bigl[(1+t)C-2(1-t)\bigr],
\]
which is not identically zero and vanishes at $\rho^{-1}$.  Divisor-form
Hurwitz forces $u_j\to\rho^{-1}$, i.e. $\tau_j-\rho_j\to0$.
Conversely, \eqref{eq:typeIIIbudget} makes every collision residue tend to
zero.  If $\rho=1$, the limiting pole $\rho^{-1}=1$ lies on $\partial\D$ and
is separated from $u=\theta^{-1}$.  If $\lambda>0$, the limiting form of
\eqref{eq:typeIIIbudget} forces $\theta-\rho>0$, so again
$u<\rho^{-1}$.  In either case meromorphic Hurwitz retains the output zero;
the same cleared polynomial $H_{\mathrm I}$, the coprime-pole check
\eqref{eq:coprimepole}, and Lemma~\ref{lem:reversal} then rule out any
further loss mode.

Finally the monic denominator has only the reciprocal pair
$\rho_j,\rho_j^{-1}$ off the unit circle.  Its noncyclotomic irreducible
factor is therefore Salem, or reciprocal quadratic Pisot.
\end{proof}

\begin{corollary}[Closed Type III strata]\label{cor:typeIIIstrata}
The limit belongs to $\calS\cup\calT$ if either
\begin{enumerate}[label=\textup{(\roman*)}]
\item $\liminf_j\lambda_j>0$, or
\item the degrees $[\Q(\rho_j):\Q]$ are bounded.
\end{enumerate}
Consequently a genuinely bad Type III sequence must satisfy
\begin{equation}\label{eq:twin}
 \rho_j\in\calT,\quad [\Q(\rho_j):\Q]\to\infty,\quad
 \rho_j,\tau_j\to\theta,\quad \lambda_j\to0.
\end{equation}
\end{corollary}

\begin{proof}
On a compact $\tau$-range, \eqref{eq:typeIIIbudget} bounds $\lambda_j$ by a
fixed multiple of $\tau_j-\rho_j$; a positive residue lower bound therefore
implies separation.  For bounded degree, $1<\rho_j<\tau_j=O(1)$ and the
Salem/Pisot property give bounded absolute Weil height.  Northcott finiteness
\cite{northcott} leaves finitely many $\rho_j$.  If collision occurs, a
constant subsequence has $\theta=\rho\in\calS\cup\calT$; otherwise apply
Theorem~\ref{thm:typeIII}.
\end{proof}

The information discarded at collision has a finite, nonzero first
derivative.

\begin{proposition}[Strict first collision tangent]\label{prop:tangent}
Assume the collision regime, pass to a subsequence for which $C_j\to C$
locally uniformly on $\D$, and suppose $\rho_j,\tau_j\to\theta>1$.  Then
\begin{equation}\label{eq:tangent}
 \frac{\lambda_j}{\tau_j-\rho_j}\longrightarrow
 L_{\theta,C}:=\frac1{2\theta}
 \left(\frac{2(\theta-1)}{\theta+1}-C(\theta^{-1})\right),
\end{equation}
and
\begin{equation}\label{eq:tangentwindow}
 0<L_{\theta,C}\le
 \frac{\theta-1}{2\theta(\theta+1)}.
\end{equation}
In particular
\[
 \lambda_j=L_{\theta,C}(\tau_j-\rho_j)
 +o(\tau_j-\rho_j).
\]
\end{proposition}

\begin{proof}
Put $u_j=\tau_j^{-1}$ and $r_j=\rho_j^{-1}$.  Solving
\eqref{eq:typeIIIdefect0} gives the exact identity
\begin{equation}\label{eq:exacttangent}
 \frac{\lambda_j}{\tau_j-\rho_j}
 =\frac{u_j(1-r_ju_j)}{2(1-u_j^2)}
 \left(\frac{2(1-u_j)}{1+u_j}-C_j(u_j)\right).
\end{equation}
Local uniform convergence and $u_j,r_j\to r:=\theta^{-1}$ yield
\eqref{eq:tangent}.

It remains to show that the limiting defect is nonzero.  Set
\[
 D_j(t)=(1-\rho_jt)(1-\rho_j^{-1}t),\qquad K_j=D_jG_j.
\]
The simple pole of $G_j$ at $r_j$ is cancelled and
\[
 K_j(r_j)=2\lambda_j(1+r_j)(1-r_j^2)\ne0.
\]
By Lemma~\ref{lem:IIIsimple}, $K_j$ has exactly one disk zero counted with
multiplicity, namely $u_j$.  On compact subsets of $\D$,
\[
 K_j\longrightarrow K(t):=(1-\theta t)(1-\theta^{-1}t)
 \bigl((1+t)C(t)-2(1-t)\bigr).
\]
The first factor gives a zero of $K$ at $r$.  On a small closed disk about
$r$ whose boundary is zero-free, the argument principle shows that this zero
has multiplicity one.  Since $1-\theta t$ already has a simple zero there,
\[
 (1+r)C(r)-2(1-r)\ne0.
\]
The bracket in \eqref{eq:exacttangent} is positive before passage to the
limit; nonvanishing makes its limit strictly positive.  Hence
$L_{\theta,C}>0$.

Finally $\lambda_j\to0$ gives $C(0)=1$, and Harnack's inequality yields
\[
 C(r)\ge\frac{1-r}{1+r}=\frac{\theta-1}{\theta+1}.
\]
Substitution in \eqref{eq:tangent} proves \eqref{eq:tangentwindow}.
\end{proof}

\begin{corollary}[Uniform first-order Type III collision]
\label{cor:IIIlinear}
Along every Type III collision with $\tau_j\to\theta>1$,
\[
 0<\liminf_j\frac{\lambda_j}{\tau_j-\rho_j}
 \le\limsup_j\frac{\lambda_j}{\tau_j-\rho_j}
 \le\frac{\theta-1}{2\theta(\theta+1)}.
\]
Thus $\lambda_j\asymp\tau_j-\rho_j$; the sublinear collision
$\lambda_j=o(\tau_j-\rho_j)$ is impossible.
\end{corollary}

\begin{proof}
The upper bound follows from \eqref{eq:typeIIIharnackbudget}.  If the lower
limit were zero, choose a subsequence on which the ratio tends to zero and
then a further locally convergent Carath\'eodory subsequence.  Proposition
\ref{prop:tangent} would give a strictly positive limit, a contradiction.
\end{proof}

The right side depends on the limiting Carath\'eodory function, not only on
$\theta$.  A useful algebraic model for the loss of an integer germ is
obtained from $A=z^2-z-1$, $A^*=1-z-z^2$, and
\[
 P_k=\frac{z^kA-A^*}{z-1},\qquad Q_k=(z-1)P_k.
\]
Direct reversal gives
\[
 Q_k^*-(1-t)P_{k+1}^*=t^k(t-1)A(t).
\]
Thus every fixed Taylor prefix stabilizes while the zero of $A$ can remain
visible only at the moving scale.  We use this identity solely as a
coefficient-compactness diagnostic; no admissibility claim for these pairs is
needed in any theorem of the paper.

\section{Type IV: residue budget and compactification}

Use throughout the Type IV real variable
\begin{equation}\label{eq:xnormalization}
 x=\sqrt z+\frac1{\sqrt z}.
\end{equation}
For a chosen Type IV witness its real quotient has the exact signed Cauchy
form
\begin{equation}\label{eq:typeIVreal}
 \mathcal R(x)=\int_{[-2,2]}\frac{d\nu(y)}{x-y}
 -\frac c{x-a}-\frac c{x+a},\qquad a>2,\quad c>0,
\end{equation}
where $\nu$ is a positive symmetric atomic measure and
\begin{equation}\label{eq:mass}
 \nu([-2,2])-2c=1.
\end{equation}
The output point $x_\tau=\sqrt\tau+\tau^{-1/2}$ lies in $(2,a)$ and obeys
\begin{equation}\label{eq:typeIVoutput}
 \mathcal R(x_\tau)=\frac2{x_\tau}.
\end{equation}

\begin{proposition}[Exact Type IV defect and universal budget]
\label{prop:IVbudget}
For a Type IV witness put
\[
 \Delta=\int_{[-2,2]}\frac{y^2}{x_\tau^2-y^2}\,d\nu(y).
\]
Then $0\le\Delta<1$ and
\begin{equation}\label{eq:IVdefect}
 \Delta+\frac{2ca^2}{a^2-x_\tau^2}=1,
 \qquad
 c=\frac{a^2-x_\tau^2}{2a^2}(1-\Delta).
\end{equation}
Consequently
\begin{equation}\label{eq:IVbudget}
 c\le\frac{a^2-x_\tau^2}{2a^2}<\frac12,\qquad
 \|\nu-c\delta_a-c\delta_{-a}\|_{\TV}=1+4c<3.
\end{equation}
Moreover
\begin{equation}\label{eq:IVmoment}
 1-2c=\Delta+\frac{2cx_\tau^2}{a^2-x_\tau^2},
 \qquad
 \int_{[-2,2]}y^2\,d\nu(y)\le x_\tau^2(1-2c).
\end{equation}
Consequently
\begin{equation}\label{eq:normalbound}
 |\mathcal R(z)|<\frac3{|\operatorname{Im}z|}\qquad(z\notin\R).
\end{equation}
\end{proposition}

\begin{proof}
For $x>2$, symmetry gives
\[
 \int\frac{d\nu(y)}{x-y}
 =\int\frac{x\,d\nu(y)}{x^2-y^2}.
\]
At $x=x_\tau<a$, the exterior pair contributes
$2cx/(a^2-x^2)$.  Multiplying the output equation by $x_\tau$
and using \eqref{eq:mass} gives exactly
\[
 2=1+2c+\Delta+\frac{2cx_\tau^2}{a^2-x_\tau^2}
  =1+\Delta+\frac{2ca^2}{a^2-x_\tau^2},
\]
which proves \eqref{eq:IVdefect}; because $c>0$, it also gives
$\Delta<1$.  Rearrangement gives the first identity in
\eqref{eq:IVmoment}.  Since
\[
 \frac{y^2}{x_\tau^2-y^2}\ge\frac{y^2}{x_\tau^2},
\]
the moment bound follows.  The total variation is
$\nu([-2,2])+2c=1+4c$, and the Cauchy-kernel bound gives
\eqref{eq:normalbound}.
\end{proof}

Consider a sequence $\tau_j\to\theta>1$ with chosen Type IV witnesses and
write $x_j=x_{\tau_j}\to x_0>2$.  By the budget, after passing to a
subsequence,
\[
 a_j\to a\in[x_0,\infty],\qquad c_j\to c\in[0,1/2],\qquad
 \nu_j\Rightarrow\nu.
\]
The Type IV ordering makes $x_j$ the unique zero in $(2,a_j)$ of
$H_j(x)=x\mathcal R_j(x)-2$, and
\begin{equation}\label{eq:onesidedsign}
 H_j(x)>0\qquad(x_j<x<a_j).
\end{equation}

\begin{theorem}[Exact Type IV compactification]\label{thm:typeIVcompact}
Every subsequential limit belongs to exactly one of the following strata.
\begin{enumerate}[label=\textup{(\roman*)}]
\item \emph{Finite separated:} $x_0<a<\infty$, and necessarily
$0<c<1/2$.
\item \emph{Zero-residue collision:} $a=x_0$ and $c_j\to0$.
\item \emph{Half-residue escape:} $a_j\to\infty$, and necessarily
\begin{equation}\label{eq:escapelimit}
 c_j\to\frac12,\qquad \nu_j\Rightarrow2\delta_0,\qquad
 \mathcal R_j\longrightarrow\frac2x
\end{equation}
locally uniformly on $\mathbb C\setminus[-2,2]$.
\end{enumerate}
\end{theorem}

\begin{proof}
The budget immediately turns $a_j-x_j\to0$ into $c_j\to0$.  Suppose now
$x_0<a<\infty$ but $c=0$.  Set
\[
 J(x)=x\int_{[-2,2]}\frac{d\nu(y)}{x-y}
 =\int_{[-2,2]}\frac{x^2}{x^2-y^2}\,d\nu(y).
\]
Its derivative is
\begin{equation}\label{eq:Jderivative}
 J'(x)=-2x\int_{[-2,2]}\frac{y^2}{(x^2-y^2)^2}\,d\nu(y)\le0.
\end{equation}
Separation lets us pass to the output
equation, giving $J(x_0)=2$, while \eqref{eq:onesidedsign} gives
$J(x)\ge2$ for $x_0<x<a$.  Monotonicity gives the reverse inequality, so
$J\equiv2$ on an interval.  Formula \eqref{eq:Jderivative} forces
$\nu$ to be supported at zero, and $J=2$ then gives $\nu=2\delta_0$,
contradicting
\eqref{eq:mass}, which would give mass one when $c=0$.  Hence $c>0$.

If $a_j\to\infty$, the exterior terms vanish at every fixed $x>2$.
The same one-sided argument gives $J(x_0)=2$ and $J(x)\ge2$ for every
$x>x_0$, hence $J\equiv2$ and $\nu=2\delta_0$.  The mass identity then gives
$c=1/2$.  These alternatives exhaust the convergent subsequences.
\end{proof}

The escape is singular at the arithmetic base point.  Write
\[
 a_j=\sqrt{\eta_j}+\eta_j^{-1/2},\qquad \eta_j>1.
\]
The disk decomposition is
\begin{equation}\label{eq:typeIVdisk}
 \Phi_j(t)=C_j(t)-
 \frac{2c_j(1-t^2)}
 {(1-\eta_jt)(1-\eta_j^{-1}t)},\qquad
 C_j(0)=1+2c_j<2.
\end{equation}
In the escape regime $\eta_j^{-1}\to0$ and
\begin{equation}\label{eq:puncturedlimit}
 \Phi_j(t)\longrightarrow\frac{2(1-t)}{1+t}
\end{equation}
locally uniformly on compact subsets of $\D\setminus\{0\}$,
whereas $\Phi_j(0)=1$ and the right side has value $2$ at zero.  The integer
germ is literally lost at the puncture.  Indeed
$\nu_j\Rightarrow2\delta_0$ in the real Cauchy chart, and the
Cauchy-to-disk transform used in deriving \eqref{eq:typeIVdisk} sends this
atom to
\[
 C_j(t)\longrightarrow\frac{2(1-t)}{1+t}
\]
locally in $\D$.  For fixed $t\ne0$, the exterior term in
\eqref{eq:typeIVdisk} tends to zero because its denominator is asymptotic to
$-\eta_jt$.  At $t=0$, however, that term equals $-2c_j\to-1$, while
$C_j(0)\to2$, leaving $\Phi_j(0)=1$.

The puncture has a coefficient-level arithmetic signature.

\begin{corollary}[Coefficient fingerprint of Type IV escape]
\label{cor:IVfingerprint}
Write
\[
 \Phi_j(t)=\sum_{n\ge0}\phi_{n,j}t^n,
 \qquad C_j(t)=\sum_{n\ge0}\gamma_{n,j}t^n.
\]
In the half-residue escape regime, for each fixed $n\ge1$,
\begin{equation}\label{eq:escape-coefficients}
 \phi_{n,j}=\gamma_{n,j}-2c_j(\eta_j^n+\eta_j^{-n}),
 \qquad
 \frac{\phi_{n,j}}{\eta_j^n}\longrightarrow-1,
 \qquad
 \frac{\phi_{n,j}}{a_j^{2n}}\longrightarrow-1.
\end{equation}
In particular every fixed nonconstant integer Taylor coefficient tends to
$-\infty$.
\end{corollary}

\begin{proof}
The coefficient of $t^n$ in
\[
 \frac{1-t^2}{(1-\eta t)(1-\eta^{-1}t)}
\]
is $\eta^n+\eta^{-n}$ for $n\ge1$.  Herglotz's coefficient bound gives
$|\gamma_{n,j}|\le2(1+2c_j)<4$.  Now use
$c_j\to1/2$, $\eta_j\to\infty$, and $a_j^2/\eta_j\to1$.
\end{proof}

\begin{theorem}[Type IV nondegenerate closure]\label{thm:typeIVclosure}
If
\begin{equation}\label{eq:typeIVseparation}
 \sup_j a_j<\infty,\qquad
 \liminf_j(a_j-x_j)>0,
\end{equation}
then $\theta\in\calS\cup\calT$.
\end{theorem}

\begin{proof}
Theorem~\ref{thm:typeIVcompact} places the limit in the finite separated
stratum and keeps $c>0$.  Formula \eqref{eq:typeIVdisk}, the bounds
$C_j(0)<2$, and bounded $\eta_j$ make every fixed Taylor coefficient bounded.
The same bound puts the disk pole $\eta_j^{-1}$ outside one fixed
neighborhood of zero.  The signed-Herglotz limit has bounded characteristic,
so Lemma~\ref{lem:compactness} gives a rational integer-germ limit
$\Phi=A/B$.  Output--pole separation in \eqref{eq:typeIVseparation} and
the local nonsingularity of $s\mapsto\sqrt s+s^{-1/2}$ separate the disk
zero $u_j=\tau_j^{-1}$ uniformly from the pole $p_j=\eta_j^{-1}$.  On a
pole-avoiding exhaustion, the argument principle retains the one disk zero
and Hurwitz excludes any additional one.  The first coprimality check in
\eqref{eq:coprimepole} excludes a zero created at a reduced denominator
pole.  Thus $H_{\mathrm I}$ has the same unique disk zero and satisfies
$H_{\mathrm I}(0)=-1$.
Lemma~\ref{lem:reversal} finishes the proof.
\end{proof}

\begin{theorem}[Sublinear Type IV collision closure]
\label{thm:IVsublinear}
Let $\tau_j\to\theta>1$ be Type IV outputs with chosen witnesses, and put
\[
 x_j=\sqrt{\tau_j}+\tau_j^{-1/2},\qquad
 x_0=\sqrt\theta+\theta^{-1/2}.
\]
If
\begin{equation}\label{eq:IVsublinear}
 a_j-x_j\longrightarrow0,
 \qquad \frac{c_j}{a_j-x_j}\longrightarrow0,
\end{equation}
then $\theta\in\calS\cup\calT$.
\end{theorem}

\begin{proof}
Write $a_j=\sqrt{\eta_j}+\eta_j^{-1/2}$ with $\eta_j>1$.  The defining
map is strictly increasing on $(1,\infty)$, so collision gives
$\eta_j\to\theta$, while \eqref{eq:IVbudget} gives $c_j\to0$.  After
extraction, $\nu_j\Rightarrow\nu$ with $\nu([-2,2])=1$.  Put
\[
 \mathcal C_j(x)=\int_{[-2,2]}\frac{d\nu_j(y)}{x-y},
 \qquad
 \mathcal C(x)=\int_{[-2,2]}\frac{d\nu(y)}{x-y}.
\]
At the output,
\[
 \mathcal C_j(x_j)=\frac2{x_j}
 -\frac{2x_j}{a_j+x_j}\frac{c_j}{a_j-x_j}.
\]
The second term tends to zero, and the Cauchy kernels converge uniformly for
$x$ near $x_0>2$.  Hence $\mathcal C(x_0)=2/x_0$.

Let $C_j$ and $C$ be the disk Carath\'eodory transforms of $\nu_j$ and
$\nu$.  Then $C_j\to C$ locally uniformly in $\D$.  By
\eqref{eq:escape-coefficients}, with $c_j\to0$ and bounded $\eta_j$, every
fixed Taylor coefficient of the exterior term in \eqref{eq:typeIVdisk}
tends to zero.  Since $\Phi_j\in\Z[[t]]$, every Taylor coefficient of $C$
is an integer.  The Carath\'eodory function $C$ has bounded characteristic,
so Theorem~\ref{thm:cantor} gives
$C=A/B\in\Q(t)$ with $A,B\in\Z[t]$ and $B(0)=1$.

Put $u_j=\tau_j^{-1}$, $u_0=\theta^{-1}$, and
\[
 G_j(t)=(1+t)\Phi_j(t)-2(1-t),\qquad
 G(t)=(1+t)C(t)-2(1-t).
\]
The moving pole $p_j=\eta_j^{-1}$ also tends to $u_0$.  On every compact
subset of $\D\setminus\{u_0\}$ the exterior term tends uniformly to zero,
so $G_j\to G$ there.  Each $G_j$ has no disk zero other than $u_j$.
Because $G(0)=-1$, the zero-free form of Hurwitz shows that $G$ has no zero
in $\D\setminus\{u_0\}$.

For $t\in(0,1)$ set $x(t)=t^{-1/2}+t^{1/2}$.  The real and disk transforms
are related by
\[
 C(t)=\frac{1-t}{\sqrt t}\,\mathcal C(x(t)),
\]
and hence
\[
 G(t)=(1-t)\bigl(J(x(t))-2\bigr),\qquad
 J(x)=x\mathcal C(x)
 =\int_{[-2,2]}\frac{x^2}{x^2-y^2}\,d\nu(y).
\]
The output identity gives $G(u_0)=0$.  Moreover
\[
 J'(x_0)=-2x_0\int_{[-2,2]}
 \frac{y^2}{(x_0^2-y^2)^2}\,d\nu(y)<0.
\]
Equality would force $\nu=\delta_0$, contrary to $J(x_0)=2$.  Since
$x'(u_0)\ne0$ and $u_0\ne1$, the zero of $G$ at $u_0$ is simple.

The reduced denominator $B$ is zero-free in $\D$, since $C$ is holomorphic
there.  Therefore
\[
 H=(1+t)A-2(1-t)B
\]
has exactly the one simple disk zero $u_0$ and satisfies $H(0)=-1$.
Lemma~\ref{lem:reversal} yields $\theta\in\calS\cup\calT$.
\end{proof}

\begin{corollary}[Residue-window closure and endpoint rates]
\label{cor:IVrates}
For every $0<\varepsilon\le1/4$, Type IV closure holds on the stratum
\begin{equation}\label{eq:residuewindow}
 \varepsilon\le c_j\le\frac12-\varepsilon.
\end{equation}
Every individual witness satisfies the quantitative separation
\begin{equation}\label{eq:IVquantgap}
 a-x_\tau\ge
 x_\tau\left((1-2c)^{-1/2}-1\right).
\end{equation}
For any sequence with $x_j\to x_0>2$,
\[
 c_j\to0\iff a_j-x_j\to0,
 \qquad
 c_j\to\frac12\iff a_j\to\infty.
\]
If $a_j-x_j\to0$ and $\nu_j\Rightarrow\nu$, put
\[
 \Delta_0=\int_{[-2,2]}\frac{y^2}{x_0^2-y^2}\,d\nu(y).
\]
Then the exact collision slope is
\begin{equation}\label{eq:collisionrate}
 \frac{c_j}{a_j-x_j}\longrightarrow\frac{1-\Delta_0}{x_0}
 \le\frac1{x_0}.
\end{equation}
If $x_0>2\sqrt2$, it also obeys
\begin{equation}\label{eq:collisionlower}
 \frac{x_0^2-8}{x_0(x_0^2-4)}
 \le\lim_j\frac{c_j}{a_j-x_j}\le\frac1{x_0}.
\end{equation}
Finally, if $\theta\notin\calS\cup\calT$, every collision sequence satisfies
\begin{equation}\label{eq:badcollisionrate}
 0<\liminf_j\frac{c_j}{a_j-x_j}
 \le\limsup_j\frac{c_j}{a_j-x_j}\le\frac1{x_0}.
\end{equation}
The only other endpoint is the forced escape
\eqref{eq:escapelimit}--\eqref{eq:puncturedlimit}.
\end{corollary}

\begin{proof}
The budget gives $x_\tau/a\le\sqrt{1-2c}$, which is
\eqref{eq:IVquantgap}.  It also shows that collision implies $c_j\to0$.
Conversely, if $c_j\to0$ but the gaps do not tend to zero, an unbounded
pole subsequence would force $c_j\to1/2$, while a bounded-pole subsequence
would have a finite separated limit with $c>0$.  Both contradict
Theorem~\ref{thm:typeIVcompact}.  The equivalence at $1/2$ follows by the
same compactification argument.

The two bounds in \eqref{eq:residuewindow} now exclude collision and escape,
respectively, so Theorem~\ref{thm:typeIVclosure} applies.  In a collision,
let $\Delta_j$ be the defect in Proposition~\ref{prop:IVbudget}.  Weak
convergence and $x_j\to x_0>2$ give $\Delta_j\to\Delta_0$, while
\eqref{eq:IVdefect} gives
\[
 \frac{c_j}{a_j-x_j}
 =\frac{a_j+x_j}{2a_j^2}(1-\Delta_j).
\]
This proves \eqref{eq:collisionrate}.  In collision the limiting measure has
mass one; hence
\[
 \Delta_0\le\frac4{x_0^2-4}.
\]
This yields \eqref{eq:collisionlower} when $x_0>2\sqrt2$.

If a bad collision had zero lower slope, a subsequence would satisfy the
hypotheses of Theorem~\ref{thm:IVsublinear}, contradicting
$\theta\notin\calS\cup\calT$.  The upper limit follows directly by dividing
\eqref{eq:IVbudget} by $a_j-x_j$.
\end{proof}

\section{A witness-independent normal-height constraint}

Endpoint localization can be sharpened independently of the choice of a
McKee--Smyth witness.  Let
$R\in\Z[z]$ be a Salem minimal polynomial of degree $2s$, and write uniquely
\begin{equation}\label{eq:tracepoly}
 R(z)=z^s r(z+z^{-1}),\qquad r\in\Z[X].
\end{equation}
The polynomial $r$ is monic irreducible of degree $s$: monicity and
integrality follow from reciprocal descent, while a proper factor of $r$
would lift to a proper factor of the Salem minimal polynomial $R$.  Its roots
are the exterior trace $y=\tau+\tau^{-1}>2$ and $s-1$ simple roots in
$(-2,2)$.  Since a Salem minimal polynomial has no cyclotomic factor, zero is
not a root of $r$.  Thus the roots of $f=Xr$ are
\[
 \alpha_1<\cdots<\alpha_s\quad(-2<\alpha_i<2),\qquad
 y>2,
\]
where one of the compact roots $\alpha_i$ is the added root $X=0$.

\begin{definition}[Normal height]
A strict integral normal interlacer is a polynomial
$b\in\Z[X]$ of degree $s-1$, with negative leading coefficient and one
simple zero in each compact gap $(\alpha_i,\alpha_{i+1})$.  Set
\begin{equation}\label{eq:kappa}
 \kappa(R)=\min_b\bigl(-\operatorname{lc}b\bigr).
\end{equation}
\end{definition}

The minimum is well defined.  Indeed, choose rational points
$\xi_i\in(\alpha_i,\alpha_{i+1})$ and clear denominators in
$-\prod_{i=1}^{s-1}(X-\xi_i)$.  The possible values of
$-\operatorname{lc}b$ form a nonempty subset of $\Z_{>0}$.  A minimizing
interlacer is primitive, since division by a nontrivial content would lower
its leading height while preserving its zeros.

\begin{theorem}[Bounded normal-height closure]\label{thm:heightclosure}
Let $R_j$ be Salem minimal polynomials with roots $\tau_j\to\theta>1$.  If
$\kappa(R_j)$ is bounded, then $\theta\in\calS\cup\calT$.
\end{theorem}

\begin{proof}
Pass to a subsequence with a fixed leading coefficient
$c=-\kappa<0$ for a minimizing interlacer $b_j$.  If
$b_j=c\prod_{k=1}^{s_j-1}(X-\xi_{k,j})$ with
$\alpha_{k,j}<\xi_{k,j}<\alpha_{k+1,j}$, direct sign counting gives positive
residues at every $\alpha_{i,j}$ and a negative residue at $y_j$.  Since
$b_j/f_j=O(X^{-2})$ at infinity, the sum of all residues is zero.  We may
therefore write uniquely
\begin{equation}\label{eq:signedresidues}
 \frac{b_j(X)}{f_j(X)}=
 \sum_{i=1}^{s_j}\frac{\lambda_{i,j}}{X-\alpha_{i,j}}
 -\frac{\Lambda_j}{X-y_j},\qquad
 \lambda_{i,j},\Lambda_j>0,\quad
 \sum_i\lambda_{i,j}=\Lambda_j.
\end{equation}
Because $f_j$ is monic integral and $\deg f_j-\deg b_j=2$,
\[
 \frac{b_j(X)}{f_j(X)}\in X^{-2}\Z[[X^{-1}]].
\]
Thus all signed power moments are integers.  For the integer Chebyshev trace
polynomials
\[
 C_0=2,\quad C_1=X,\quad C_{n+1}=XC_n-C_{n-1},
\]
the signed moments
\begin{equation}\label{eq:integercheb}
 a_{j,n}=\sum_i\lambda_{i,j}C_n(\alpha_{i,j})
 -\Lambda_jC_n(y_j)
\end{equation}
the signed moments are therefore integers as well.  For fixed $n$,
\begin{equation}\label{eq:momentbound}
 |a_{j,n}|\le\Lambda_j\bigl(2+C_n(y_j)\bigr),
\end{equation}
which is uniformly bounded by \eqref{eq:massbound} below and
$y_j\to\theta+\theta^{-1}$.

Let $\bar\alpha_j=\Lambda_j^{-1}\sum_i\lambda_{i,j}\alpha_{i,j}\in[-2,2]$.
The first moment is
$c=\Lambda_j(\bar\alpha_j-y_j)$, so
\begin{equation}\label{eq:massbound}
 \frac{|c|}{y_j+2}\le\Lambda_j\le\frac{|c|}{y_j-2}.
\end{equation}
Because $\theta>1$, the masses are bounded above and away from zero.  Lift
$\alpha_{i,j}=\zeta_{i,j}+\zeta_{i,j}^{-1}$ to the unit circle and form the
symmetric positive measure
\[
 \mu_j=\frac12\sum_i\lambda_{i,j}
 (\delta_{\zeta_{i,j}}+\delta_{\zeta_{i,j}^{-1}}).
\]
After weak-* and diagonal extraction, $\Lambda_j\to\Lambda>0$ and every
integer $a_{j,n}$ stabilizes to $a_n$.  The Herglotz function
\[
 U(t)=2\int_{\T}\frac{d\mu(\zeta)}{1-t\zeta}
\]
has positive real part.  Hence $(U-1)/(U+1)\in H^\infty$, so $U$ has
bounded characteristic.  The function
\begin{equation}\label{eq:Aseries}
 \mathcal A(t)=\sum_{n\ge0}a_nt^n
 =U(t)-\frac{\Lambda}{1-\theta t}
       -\frac{\Lambda}{1-\theta^{-1}t}
\end{equation}
extends meromorphically with bounded characteristic to $\D$.
Theorem~\ref{thm:cantor} makes $\mathcal A$ rational.  The term
$-\Lambda/(1-\theta t)$ gives a genuine pole at $\theta^{-1}$, because $U$
and the other rational term are holomorphic there and $\Lambda>0$.  It is the
only pole in $\D$.  If $A/B$ is the reduced integral form with $B(0)=1$,
then
\[
 B^\vee(z)=z^{\deg B}B(z^{-1})
\]
is monic integral, has the unique root $\theta$ outside the closed unit disk,
and has every other root in that disk.  The irreducible factor containing
$\theta$ makes it Pisot or Salem.
\end{proof}

There is also a uniform linear upper bound, so a hypothetical bad sequence
is confined to a linear-height cusp rather than an arbitrary one.

\begin{theorem}[Linear normal-height bound]\label{thm:linearheight}
For a Salem minimal polynomial of degree $2s$ and Salem root $\tau$,
\begin{equation}\label{eq:linearheight}
 \kappa(R)<(16\tau^2+48)(s-1).
\end{equation}
\end{theorem}

\begin{proof}
Let $C_m$ be as above, put
\[
 H_m=2+C_m,\qquad h_m=(2-X)H_m,\qquad
 D_m=(y-2)H_m(y),
\]
where $C_m(y)=\tau^m+\tau^{-m}$.  Thus $D_m\to\infty$, so the least
$m\ge0$ with $D_m>16(s-1)$ exists; recall also that $s\ge2$ for a Salem
polynomial.  Define the integral
remainder
\begin{equation}\label{eq:bm}
 b_m=\operatorname{rem}_{r}\bigl[Xr'(X)h_m(X)\bigr].
\end{equation}
Because $r$ is monic, Euclidean division preserves integrality, so
$b_m\in\Z[X]$ and $\deg b_m<s$.
For a compact conjugate $\beta\in(-2,2)$, one has
$0<H_m(\beta)\le4$; if $H_m(\beta)=0$, irreducibility would force $r$ to
divide $H_m$, impossible at $y$.  Hence
\[
 0<h_m(\beta)<16,\qquad |\beta h_m(\beta)|<32,\qquad h_m(y)=-D_m.
\]
At each root $\gamma$ of $r$,
$b_m(\gamma)/(\gamma r'(\gamma))=h_m(\gamma)$.  The residue at every compact
pole is positive and the exterior residue is negative.  Lagrange
interpolation at the roots of the monic polynomial $r$ gives the two exact
trace identities
\begin{align}
 \frac{b_m(0)}{r(0)}
 &=-\operatorname{Tr}_{\Q(y)/\Q}h_m
   >D_m-16(s-1)>0,\label{eq:zeropoletrace}\\
 [X^{s-1}]b_m
 &=\operatorname{Tr}_{\Q(y)/\Q}(Xh_m),\qquad
 \left|[X^{s-1}]b_m+yD_m\right|<32(s-1).
 \label{eq:lctrace}
\end{align}
Since $y>2$ and $D_m>16(s-1)$, the second estimate in
\eqref{eq:lctrace} also gives
$[X^{s-1}]b_m<-yD_m+32(s-1)<0$.  Thus $\deg b_m=s-1$ and this
coefficient is $\operatorname{lc}b_m$.
The first identity is the residue at the added zero pole $X=0$.  Thus all
compact residues of $b_m/f$ are positive and the exterior residue is
negative.  The function crosses from one infinite sign to the other in every
compact gap; since $\deg b_m=s-1$, it has exactly one simple zero in every
such gap and is a strict normal interlacer.

If $m\ge1$, minimality gives $D_{m-1}\le16(s-1)$.  Since
$\tau C_{m-1}(y)-C_m(y)=\tau^{-m}(\tau^2-1)>0$ and $2<2\tau$,
one has
\[
 D_m<\tau D_{m-1}\le16\tau(s-1).
\]
For $m=0$, use $y-2=(\tau-1)^2/\tau<\tau$ and $s\ge2$ to get
$D_0=4(y-2)<4\tau\le16\tau(s-1)$.  Combining this with
\eqref{eq:lctrace} and $y\tau=\tau^2+1$ gives
\[
 \kappa(R)\le-\operatorname{lc}b_m
 <(16y\tau+32)(s-1)=(16\tau^2+48)(s-1).
\]
\end{proof}

\begin{corollary}[Exact computation at fixed input]
\label{cor:kappacomputable}
For every fixed Salem minimal polynomial $R$, the integer $\kappa(R)$ is
effectively computable by a finite exact search.
\end{corollary}

\begin{proof}
Put $n=s-1$ and $B_R=(16\tau^2+48)n$.  Theorem
\ref{thm:linearheight} reduces the leading height to
$1\le K\le\lceil B_R\rceil-1$; the ceiling is decidable by exact real-root
isolation.  If
\[
 b=-K\prod_{i=1}^{n}(X-\xi_i)
\]
is a strict normal interlacer, then $|\xi_i|<2$, so Vieta's formulas give
\[
 \bigl|[X^{n-q}]b\bigr|\le K\binom{n}{q}2^q
 \qquad(0\le q\le n).
\]
Only finitely many integral candidates occur.  Sturm sequences and exact
root isolation decide whether a candidate has one simple zero in every
prescribed compact gap.  Searching increasing $K$ returns $\kappa(R)$.
\end{proof}

\section{Aggregate endpoint localization}

Recall the four output classes defined in Section~2, always with one witness
selected for each term.  Type I is closed by
Proposition~\ref{prop:typeI}, and Type II is closed by
Theorem~\ref{thm:typeII}.

\begin{theorem}[Three-cusp localization]\label{thm:aggregate}
Suppose $\tau_j\in\calT$ and $\tau_j\to\theta>1$ with
$\theta\notin\calS\cup\calT$.  Choose one McKee--Smyth witness for every
$j$, pass to a subsequence of fixed type, and then, if necessary, to a
further subsequence on which the pole, residue, and compact-measure
parameters converge.  Exactly one of the following descriptions holds for
the resulting subsequence.  In the Type IV alternatives, put
\[
 x_j=\sqrt{\tau_j}+\tau_j^{-1/2},\qquad
 x_0=\sqrt\theta+\theta^{-1/2}.
\]
\begin{enumerate}[label=\textup{(\Roman*)}]
\item \emph{Type III Salem-twin collision:}
\[
 \rho_j\in\calT,\quad [\Q(\rho_j):\Q]\to\infty,\quad
 \rho_j,\tau_j\to\theta,\quad \tau_j-\rho_j\to0,\quad
 \lambda_j\to0
\]
\[
 0<\liminf_j\frac{\lambda_j}{\tau_j-\rho_j}
 \le\limsup_j\frac{\lambda_j}{\tau_j-\rho_j}
 \le\frac{\theta-1}{2\theta(\theta+1)}.
\]
\item \emph{Type IV zero-residue collision:}
\[
 a_j-x_j\to0,\quad c_j\to0,\quad
 0<\liminf_j\frac{c_j}{a_j-x_j}
 \le\limsup_j\frac{c_j}{a_j-x_j}\le\frac1{x_0}.
\]
\item \emph{Type IV half-residue escape:}
\[
 a_j\to\infty,\quad c_j\to\frac12,\quad
 \nu_j\Rightarrow2\delta_0,\quad \mathcal R_j\to2/x
 \text{ locally uniformly on }\mathbb C\setminus[-2,2],
\]
with loss of the integer germ at $t=0$ as in \eqref{eq:puncturedlimit}.
\end{enumerate}
In every case, if $R_j$ is the Salem minimal polynomial of $\tau_j$ and
$2s_j=\deg R_j$, then
\begin{equation}\label{eq:aggregateheight}
 s_j\to\infty,\qquad \kappa(R_j)\to\infty,\qquad
 \kappa(R_j)<(16\tau_j^2+48)(s_j-1).
\end{equation}
\end{theorem}

\begin{proof}
There are only four witness types.  A fixed-type subsequence cannot be Type I
or II by their closure results.  Theorem~\ref{thm:typeIII} and
Corollaries~\ref{cor:typeIIIstrata} and \ref{cor:IIIlinear} give (I).
Theorems~\ref{thm:typeIVcompact}, \ref{thm:typeIVclosure}, and
\ref{thm:IVsublinear}, together with Corollary~\ref{cor:IVrates}, give (II)
or (III).  These are alternatives only within the resulting
parameter-convergent subsequence; a different witness choice or subsequence
may produce a different local chart.

If $\kappa(R_j)$ had a bounded subsequence, Theorem~\ref{thm:heightclosure}
would put $\theta$ in $\calS\cup\calT$.  Thus the height tends to infinity
along every subsequence.  If $s_j$ were bounded, then
$h(\tau_j)=\log\tau_j/(2s_j)$ and the degrees would both be bounded;
Northcott finiteness would leave a constant subsequence, contradicting the
bad limit.  Hence $s_j\to\infty$.  The final inequality is
Theorem~\ref{thm:linearheight}; since $\tau_j\to\theta$, it also says
$\kappa(R_j)=O_\theta(s_j)$.
\end{proof}

\section{What remains}

Theorem~\ref{thm:aggregate} converts an amorphous closure problem into three
specific endpoint problems, but it does not eliminate them.  A complete
proof of interior closure must eliminate all three, either through a common
arithmetic rigidity principle or through separate mechanisms adapted to the
three cases:
\begin{itemize}
\item cross-degree rigidity of the strictly positive Type III tangent
\eqref{eq:tangent} in unbounded-degree Salem-twin collisions;
\item recovery of the cancelled pole in the strictly first-order Type IV
zero-residue collision; and
\item recovery of an arithmetic germ through the Type IV escape puncture,
despite the coefficient growth in \eqref{eq:escape-coefficients}.
\end{itemize}
Coefficient stabilization alone does not exclude the first or third regime,
as the explicit collision identity and \eqref{eq:puncturedlimit} demonstrate.

There is a precise interface with Mahler-function arithmetic that may be
useful after a terminal rescaling, but the rescaling itself remains missing.

\begin{proposition}[Adelic Mahler pole interface]\label{prop:Mahlerinterface}
Let $F\in\Q[[t]]$ be $k$-Mahler for some fixed $k\ge2$.  Assume that its
$p$-adic Taylor radius is at least one for every rational prime $p$, and that
its complex germ is meromorphic in $\D$ with exactly one pole
$\rho\in(0,1)$.  Then $\rho^{-1}$ is a Pisot or Salem number.
\end{proposition}

\begin{proof}
Let $d(t)$ be the normalized Mahler denominator of $F$ in the sense of
\cite[Def.~3.9]{abs}.  Galois invariance and normalization give
$d\in\Q[t]$.  The denominator identity shows that $d(\rho)=0$: otherwise
all iterates $F(t^{k^i})$ are holomorphic at $\rho$, contradicting the pole.
The finite-place part of \cite[Prop.~6.4]{abs} and the assumed radii show
that $\rho^{-1}$ is an algebraic integer.  The same local denominator
argument at the complex embeddings shows that every other conjugate of
$\rho^{-1}$ has modulus at most one, since a second conjugate of $\rho$
strictly inside $\D$ would be a second pole.  If all inequalities are strict
the number is Pisot.  Otherwise irreducibility and a unit-circle conjugate
force the minimal polynomial to be reciprocal, leaving precisely the
reciprocal pair off the unit circle; the number is Salem.
\end{proof}

Proposition~\ref{prop:Mahlerinterface} is conditional at the Salem--Boyd
cusps: this paper constructs neither a fixed-base Mahler equation nor the
required all-place analytic germ in any of the three terminal regimes.  It
therefore supplies a target criterion, not an additional closure claim.

All conclusions in this paper are relative to $\theta>1$.  Even a resolution
of the three endpoints would establish only
\[
 \overline{\calT}\cap(1,\infty)\subseteq\calS\cup\calT.
\]
It would not exclude accumulation at $1$, and it would not prove the stronger
derived-set identity $(\calS\cup\calT)'=\calS$.

\section*{Verification note}

The residue-defect identities, the Type II cleared identity, the two Weyl
complements in Proposition~\ref{prop:weylfailure}, the collision changes of
variables, and the escape coefficients follow by direct expansion from the
displayed formulas.  No theorem in this paper depends on numerical or
computer-algebra evidence.  Compactness, rationality, zero retention, and
Northcott finiteness are proved or cited in the text.

\begin{thebibliography}{99}

\bibitem{abs}
B. Adamczewski, J. P. Bell and D. Smertnig,
\emph{A height gap theorem for coefficients of Mahler functions},
J. Eur. Math. Soc. (JEMS) \textbf{25} (2023), no.~7, 2525--2571,
\href{https://doi.org/10.4171/JEMS/1244}{doi:10.4171/JEMS/1244}.

\bibitem{amara25}
M. Amara,
\emph{Nombres de Pisot, nombres de Salem et la conjecture de Lehmer},
arXiv:2509.21402v1 [math.NT] (2025),
\href{https://arxiv.org/abs/2509.21402}{arXiv:2509.21402}.

\bibitem{boyd77}
D. W. Boyd,
\emph{Small Salem numbers},
Duke Math. J. \textbf{44} (1977), no.~2, 315--328,
\href{https://doi.org/10.1215/S0012-7094-77-04413-1}
{doi:10.1215/S0012-7094-77-04413-1}.

\bibitem{boyd78}
D. W. Boyd,
\emph{Pisot and Salem numbers in intervals of the real line},
Math. Comp. \textbf{32} (1978), no.~144, 1244--1260,
\href{https://doi.org/10.1090/S0025-5718-1978-0491587-8}
{doi:10.1090/S0025-5718-1978-0491587-8}.

\bibitem{cantor}
D. G. Cantor,
\emph{Power series with integral coefficients},
Bull. Amer. Math. Soc. \textbf{69} (1963), 362--366,
\href{https://doi.org/10.1090/S0002-9904-1963-10923-4}
{doi:10.1090/S0002-9904-1963-10923-4}.

\bibitem{mckeesmyth}
J. McKee and C. Smyth,
\emph{Salem numbers and Pisot numbers via interlacing},
Canad. J. Math. \textbf{64} (2012), no.~2, 345--367,
\href{https://doi.org/10.4153/CJM-2011-051-2}
{doi:10.4153/CJM-2011-051-2}.

\bibitem{northcott}
D. G. Northcott,
\emph{An inequality in the theory of arithmetic on algebraic varieties},
Proc. Cambridge Philos. Soc. \textbf{45} (1949), no.~4, 502--509,
\href{https://doi.org/10.1017/S0305004100025202}
{doi:10.1017/S0305004100025202}.

\end{thebibliography}

\end{document}
